% ===================== PREAMBULE COMMUN (à copier VERBATIM en tête de CHAQUE .tex) =====================
\documentclass[11pt,a4paper]{article}
\usepackage[utf8]{inputenc}
\usepackage[T1]{fontenc}
\usepackage{lmodern}
\usepackage[french]{babel}
\usepackage{amsmath,amssymb,mathtools}
\usepackage{icomma}
\usepackage{xcolor}
\usepackage{colortbl}
\usepackage{tikz}
\usepackage{pgfplots}
\usepackage[most]{tcolorbox}
\usepackage{enumitem}
\usepackage{array}
\usepackage{booktabs}
\usepackage{multicol}
\usepackage{fancyhdr}
\usepackage{caption}
\captionsetup{labelformat=empty,justification=centering,font=small}
\usepackage[hidelinks]{hyperref}
\usetikzlibrary{arrows.meta,positioning,calc,angles,quotes,patterns,backgrounds,shapes.geometric}
\pgfplotsset{compat=1.18}
\usepackage[a4paper,margin=2.2cm,headheight=26pt,headsep=10pt]{geometry}

\definecolor{primary}{RGB}{219,54,77}
\definecolor{primarydark}{RGB}{150,28,46}
\definecolor{primarylight}{RGB}{250,224,228}
\definecolor{defcol}{RGB}{0,114,178}
\definecolor{thmcol}{RGB}{0,140,120}
\definecolor{excol}{RGB}{0,135,90}
\definecolor{methcol}{RGB}{90,90,100}
\definecolor{courbe}{RGB}{40,90,200}
\definecolor{courbeB}{RGB}{210,60,40}
\definecolor{courbeC}{RGB}{30,150,80}
\definecolor{axiscol}{RGB}{60,60,60}
\definecolor{gridcol}{RGB}{200,205,215}

\tcbset{boxbase/.style={enhanced,breakable,boxrule=0.9pt,arc=2.5pt,
   left=3mm,right=3mm,top=2.3mm,bottom=2.3mm,
   fonttitle=\bfseries,coltitle=white,toptitle=0.6mm,bottomtitle=0.6mm}}
\newtcolorbox{methode}[1][]{boxbase,colback=methcol!7,colframe=methcol,sharp corners,
  title={\textsc{Méthode}\ifx\relax#1\relax\relax\else\textmd{ --- \itshape #1}\fi}}
\newtcolorbox{exemple}[1][]{boxbase,colback=excol!5,colframe=excol,
  title={\textsc{Exemple}\ifx\relax#1\relax\relax\else\textmd{ : \itshape #1}\fi}}
\newtcolorbox{idee}[1][]{boxbase,colback=defcol!6,colframe=defcol,
  title={\textsc{L'essentiel}\ifx\relax#1\relax\relax\else\textmd{ : \itshape #1}\fi}}
\newtcolorbox{piege}[1][]{boxbase,colback=primary!7,colframe=primary,
  title={\textsc{Piège}\ifx\relax#1\relax\relax\else\textmd{ : \itshape #1}\fi}}

\pgfplotsset{repere/.style={axis lines=middle,
    axis line style={-{Stealth[length=2.2mm]},axiscol,line width=0.8pt},
    label style={font=\footnotesize},tick label style={font=\scriptsize},
    every axis x label/.style={at={(current axis.right of origin)},anchor=north west},
    every axis y label/.style={at={(current axis.above origin)},anchor=south east},
    xlabel={$x$},ylabel={$y$},grid=both,minor tick num=0,
    major grid style={gridcol,line width=0.4pt},samples=200,clip=false}}

\pagestyle{fancy}\fancyhf{}
\fancyhead[L]{\small\textcolor{primarydark}{\textbf{Fiche 4} — Les fonctions}}
\fancyhead[R]{\small\textcolor{primarydark}{\textsc{Carnet d'entraînement}}}
\fancyfoot[C]{\small\thepage}
\renewcommand{\headrulewidth}{0.8pt}
\renewcommand{\headrule}{\hbox to\headwidth{\color{primary}\leaders\hrule height \headrulewidth\hfill}}
\usepackage{titlesec}
\titleformat{\section}{\Large\bfseries\color{primarydark}}{\thesection}{0.6em}{}
\titleformat{\subsection}{\large\bfseries\color{primary}}{\thesubsection}{0.6em}{}

\newcommand{\R}{\mathbb{R}}\newcommand{\Z}{\mathbb{Z}}\newcommand{\N}{\mathbb{N}}
\newcommand{\Df}{D_f}
\newcommand{\croit}{\textcolor{thmcol}{\Large$\nearrow$}}
\newcommand{\decroit}{\textcolor{thmcol}{\Large$\searrow$}}

\newcommand{\exercice}[1]{\medskip\noindent{\color{primarydark}\bfseries\large Exercice #1}\par\nopagebreak\smallskip}
\newcommand{\titrecarnet}[3]{%
\begin{center}\begin{tikzpicture}
\node[rectangle,rounded corners=7pt,fill=primary,text=white,minimum width=15cm,
  minimum height=1.7cm,align=center,inner sep=8pt]
  {{\LARGE\bfseries #1}\\[2pt]{\normalsize #2}};
\end{tikzpicture}\end{center}
\begin{center}\itshape\color{primarydark}#3\end{center}\medskip}

% ---- RÈGLES D'USAGE DES MACROS ----
% * \croit et \decroit s'emploient en mode TEXTE (dans une cellule de tableau),
%   JAMAIS entre $...$ (ils contiennent déjà un $\nearrow$).
% * \titrecarnet{Titre}{Sous-titre}{Ligne en italique} : bandeau de titre.
% * \exercice{1} : titre d'exercice.
% * Boîtes : \begin{idee}...\end{idee}, \begin{exemple}[titre]...\end{exemple},
%   \begin{methode}[titre]...\end{methode}, \begin{piege}[titre]...\end{piege}.
% Le corps du document commence après \begin{document}.
\begin{document}
\titrecarnet{Thème 1 — Images, antécédents \& domaine}{Corrigé — niveau Difficile}{Détail des calculs.}

\exercice{1}
Racines de $x^2-x-6$ : $\Delta=(-1)^2-4(1)(-6)=1+24=25$, d'où $x=\dfrac{1\pm 5}{2}$, soit $x=3$ et $x=-2$. Donc $x^2-x-6=(x-3)(x+2)$.
\begin{enumerate}[label=\alph*),leftmargin=1.8em,topsep=2pt]
  \item \textbf{Racine carrée} : il faut $(x-3)(x+2)\geq 0$. Un produit de ce trinôme (avec $a=1>0$) est $\geq 0$ \emph{à l'extérieur} des racines, bornes \emph{incluses} :
        \[ \boxed{\Df=\,]-\infty,-2]\ \cup\ [3,+\infty[}. \]
  \item \textbf{Logarithme} : il faut $(x-3)(x+2)>0$ (\emph{strict}), donc à l'extérieur des racines, bornes \emph{exclues} :
        \[ \boxed{D_g=\,]-\infty,-2[\ \cup\ ]3,+\infty[}. \]
        Différence : le log \emph{exclut} $-2$ et $3$ (où le trinôme vaut $0$), alors que la racine les \emph{inclut}.
\end{enumerate}

\exercice{2}
\begin{enumerate}[label=\alph*),leftmargin=1.8em,topsep=2pt]
  \item La racine \emph{cubique} (indice \emph{impair}) existe pour tout réel : par exemple $\sqrt[3]{-8}=-2$. Aucune condition.
        \[ \boxed{\Df=\R}. \]
  \item La racine \emph{carrée} impose $2x+1\geq 0\Leftrightarrow x\geq -\tfrac12$.
        \[ \boxed{D_g=\Big[-\tfrac12,+\infty\Big[}. \]
  \item Seules les racines d'indice \emph{pair} ($\sqrt{\ },\sqrt[4]{\ },\dots$) exigent un argument $\geq 0$. Les racines d'indice \emph{impair} sont définies sur $\R$.
\end{enumerate}

\begin{piege}
Écrire $\Df=\big[-\tfrac12,+\infty\big[$ pour $\sqrt[3]{2x+1}$ est l'erreur classique : on confond racine cubique et racine carrée.
\end{piege}

\exercice{3}
Condition d'existence de $\arcsin(2x-1)$ : $-1\leq 2x-1\leq 1$.
\[ -1\leq 2x-1 \Rightarrow 0\leq 2x \Rightarrow x\geq 0, \qquad 2x-1\leq 1\Rightarrow 2x\leq 2\Rightarrow x\leq 1. \]
L'intersection donne
\[ \boxed{\Df=[0,\,1]}. \]

\exercice{4}
\begin{enumerate}[label=\alph*),leftmargin=1.8em,topsep=2pt]
  \item $f\big(f(x)\big)=\dfrac{f(x)}{f(x)+1}=\dfrac{\frac{x}{x+1}}{\frac{x}{x+1}+1}$. Au dénominateur : $\dfrac{x}{x+1}+1=\dfrac{x+(x+1)}{x+1}=\dfrac{2x+1}{x+1}$. Donc
        \[ f\big(f(x)\big)=\frac{\frac{x}{x+1}}{\frac{2x+1}{x+1}}=\frac{x}{x+1}\times\frac{x+1}{2x+1}=\boxed{\dfrac{x}{2x+1}}. \]
  \item Pour que $f\circ f$ existe, il faut d'abord que $f(x)$ existe ($x\neq -1$), puis que $f$ s'applique à $f(x)$, c'est-à-dire $f(x)\neq -1$. Or $\dfrac{x}{x+1}=-1\Leftrightarrow x=-(x+1)\Leftrightarrow 2x=-1\Leftrightarrow x=-\tfrac12$. On retrouve d'ailleurs le dénominateur $2x+1$ qui s'annule en $x=-\tfrac12$. D'où
        \[ \boxed{D_{f\circ f}=\R\setminus\Big\{-1,\,-\tfrac12\Big\}}. \]
\end{enumerate}

\exercice{5}
Deux conditions à réunir :
\begin{itemize}[leftmargin=1.6em,topsep=2pt]
  \item logarithme : $x>0$ ;
  \item racine carrée \emph{au dénominateur} : $4-x^2>0$ (strict) $\Leftrightarrow -2<x<2$.
\end{itemize}
L'intersection de $x>0$ et $-2<x<2$ donne
\[ \boxed{\Df=\,]0,\,2[}. \]
Autrement dit, l'expression n'a de sens que pour $x$ strictement compris entre $0$ et $2$.
\end{document}
